\documentclass[11pt]{article}
\usepackage[margin=1in]{geometry}
\usepackage{amsmath,amssymb}
\title{Disproof of Conjecture 00000000348\\ \large (twin partial quotients via divisor function)}
\date{October 2, 2026}
\begin{document}
\maketitle

\section{Conjecture}
There exists a real number $\alpha$ whose continued-fraction partial quotients satisfy
\[
  a_n = a_{n+1} = \lfloor \tau(n) \rfloor \qquad \text{for all } n \ge 1,
\]
where $\tau(n)$ denotes the divisor function (the number of positive divisors of $n$);
moreover $\alpha$ is given by an explicit series construction.

\section{Disproof}
\textbf{Verdict: FALSE.}

The constraints imposed by consecutive indices overlap in a single partial quotient:
the condition at $n$ fixes $a_{n+1}$ to $\lfloor\tau(n)\rfloor$, while the condition at
$n+1$ fixes the same quantity $a_{n+1}$ to $\lfloor\tau(n+1)\rfloor$. Hence such an
$\alpha$ can exist only if
\[
  \lfloor\tau(n)\rfloor = \lfloor\tau(n+1)\rfloor \quad\text{for every } n \ge 1 .
\]

Take $n = 4$ and $n = 5$:
\[
  \tau(4) = \lvert\{1,2,4\}\rvert = 3, \qquad
  \tau(5) = \lvert\{1,5\}\rvert = 2 .
\]
The condition at $n=4$ forces $a_4 = a_5 = 3$; the condition at $n=5$ forces
$a_5 = a_6 = 2$. Thus $a_5 = 3$ and $a_5 = 2$ simultaneously, a contradiction.

Therefore no real number $\alpha$ satisfies the stated partial-quotient pattern, and in
particular none is obtainable from an explicit series construction. The conjecture is
false.

\section{Remarks}
\begin{itemize}
\item Since $\tau(n) \in \mathbb{N}$, the floor is immaterial:
      $\lfloor\tau(n)\rfloor = \tau(n)$.
\item The conflict occurs even earlier: $n = 1,2$ give
      $a_2 = \tau(1) = 1$ and $a_2 = \tau(2) = 2$. The argument above uses the
      registered attack pair $(n,n+1) = (4,5)$.
\item In general the condition forces $\tau(n) = \tau(n+1)$ for all $n$, i.e.\ a
      constant divisor function, which is absurd since $\tau(1) = 1 \ne 2 = \tau(2)$;
      any adjacent pair with differing $\tau$ yields a disproof.
\end{itemize}

\section{Verification}
The numeric claims $\tau(4) = 3$ and $\tau(5) = 2$ are recomputed by
\texttt{reproduce.py} and formalized in Lean 4 (no Mathlib, zero axioms) in
\texttt{lean4/Main.lean}, which proves
\[
  \neg\,\exists\, a : \mathbb{N} \to \mathbb{N},\; \forall n,\; a_n = \tau(n) \wedge a_{n+1} = \tau(n).
\]

\end{document}
